\documentclass{dlproblemset}

\title{\textit{Graph Neural Networks}}
\subtitle{Lista de Exercícios}

\begin{document}

\paragraph{Instruções:} Responda às questões de múltipla escolha assinalando apenas uma alternativa. Nas questões dissertativas, apresente o raciocínio passo a passo, incluindo fórmulas e cálculos intermediários quando solicitado. Mantenha a notação utilizada em aula: $G = (V, E)$ para o grafo, $\mathbf{A}$ para a \textit{adjacency matrix}, $\mathbf{X}$ para os atributos de nó, $\mathcal{N}(v)$ para a vizinhança de $v$, $\mathbf{h}^{(k)}_v$ para o embedding de $v$ na camada $k$, $\tilde{\mathbf{A}} = \mathbf{A} + \mathbf{I}$ e $\tilde{\mathbf{D}} = \mathrm{diag}(\tilde{\mathbf{A}} \mathbf{1})$ para a normalização do GCN, e $\alpha_{vu}$ para os coeficientes de atenção do GAT. Em valores numéricos use a vírgula como separador decimal.

\subsection*{Questões de Múltipla Escolha.}

\begin{enumerate}[I)]

    \item Considere o grafo não-dirigido $G$ com $V = \{1,2,3,4\}$ e arestas $E = \{(1,2),(2,3),(2,4),(3,4)\}$. A partir dessa descrição, qual é o vetor de graus $(d_1, d_2, d_3, d_4)$ dos quatro nós?
    \begin{enumerate}[(a)]
        \item $(1,\, 3,\, 2,\, 2)$
        \item $(2,\, 4,\, 3,\, 3)$
        \item $(1,\, 2,\, 2,\, 1)$
        \item $(2,\, 3,\, 3,\, 2)$
        %a
    \end{enumerate}

    \item Uma camada de \textit{message passing} agrega as mensagens dos vizinhos por meio de uma função $\mathrm{AGG}$. Que propriedade essa função precisa satisfazer para que a camada seja equivariante à permutação dos nós do grafo?
    \begin{enumerate}[(a)]
        \item Ser uma transformação linear dos embeddings dos vizinhos.
        \item Ser indiferente à ordem em que os vizinhos são apresentados.
        \item Ser uma função monotônica crescente do grau do nó central.
        \item Ser injetiva sobre vetores individuais de atributos de nó.
        %b
    \end{enumerate}

    \item No grafo da questão I, considere uma agregação por média que inclui o próprio nó,
    \[
        \frac{1}{|\mathcal{N}(v) \cup \{v\}|}\sum_{u \in \mathcal{N}(v) \cup \{v\}} h^{(0)}_u,
    \]
    com atributos escalares iniciais $\mathbf{h}^{(0)} = (1,\, 2,\, 3,\, 4)^\top$. Qual é o valor agregado para o nó $2$?
    \begin{enumerate}[(a)]
        \item $2{,}5$
        \item $3{,}0$
        \item $4{,}5$
        \item $2{,}0$
        %a
    \end{enumerate}

    \item Na forma matricial do GCN, $\mathbf{H}^{(k)} = \sigma(\tilde{\mathbf{D}}^{-1/2} \tilde{\mathbf{A}} \tilde{\mathbf{D}}^{-1/2} \mathbf{H}^{(k-1)} \mathbf{W}^{(k)})$, o fator $\tilde{\mathbf{D}}^{-1/2} \tilde{\mathbf{A}} \tilde{\mathbf{D}}^{-1/2}$ é calculado uma única vez e reaproveitado a cada camada e a cada época. Qual característica do problema justifica esse reaproveitamento?
    \begin{enumerate}[(a)]
        \item Esse fator não contém parâmetros aprendíveis e depende somente da estrutura fixa do grafo.
        \item Esse fator é recalculado por \textit{backpropagation} mas converge nas primeiras épocas.
        \item Esse fator depende de $\mathbf{W}^{(k)}$, que permanece congelado durante todo o treino.
        \item Esse fator muda a cada camada porque incorpora a não-linearidade $\sigma$.
        %a
    \end{enumerate}

    \item Considere a aresta entre os nós $2$ e $3$ do grafo da questão I. No GCN, o peso aplicado à mensagem que viaja por essa aresta é $1/\sqrt{\tilde d_v \tilde d_u}$, onde $\tilde d$ é o grau no grafo com auto-laços. Qual é, aproximadamente, esse peso?
    \begin{enumerate}[(a)]
        \item $0{,}41$
        \item $0{,}29$
        \item $0{,}50$
        \item $0{,}33$
        %b
    \end{enumerate}

    \item No GCN clássico, dois nós com a mesma estrutura local e os mesmos atributos recebem embeddings idênticos após a propagação. No grafo da questão I, com atributos escalares iniciais iguais para todos os nós, qual par de nós necessariamente termina com o mesmo embedding?
    \begin{enumerate}[(a)]
        \item Os nós $1$ e $2$.
        \item Os nós $2$ e $3$.
        \item Os nós $3$ e $4$.
        \item Os nós $1$ e $4$.
        %c
    \end{enumerate}

    \item Um modelo treinado em um conjunto de grafos de moléculas precisa ser aplicado a moléculas inéditas no momento do teste. Entre GCN na forma matricial e GraphSAGE com \textit{neighbor sampling}, qual escolha é mais adequada e por qual razão estrutural?
    \begin{enumerate}[(a)]
        \item GCN, pois a matriz $\tilde{\mathbf{D}}^{-1/2} \tilde{\mathbf{A}} \tilde{\mathbf{D}}^{-1/2}$ generaliza diretamente para grafos não vistos.
        \item GraphSAGE, pois aprende uma função de agregação aplicável a vizinhanças de qualquer grafo.
        \item GCN, pois o reaproveitamento do fator de propagação reduz o custo em grafos novos.
        \item GraphSAGE, pois a amostragem de vizinhos elimina a necessidade de atributos de nó.
        %b
    \end{enumerate}

    \item No GAT, um nó $v$ tem três vizinhos cujos \textit{scores} de atenção não normalizados são $e_{v,u_1} = 2{,}0$, $e_{v,u_2} = 1{,}0$ e $e_{v,u_3} = 1{,}5$. Após a normalização por $\mathrm{softmax}$ sobre $\mathcal{N}(v)$, qual é, aproximadamente, o coeficiente $\alpha_{v,u_1}$?
    \begin{enumerate}[(a)]
        \item $0{,}33$
        \item $0{,}19$
        \item $0{,}51$
        \item $0{,}67$
        %c
    \end{enumerate}

    \item Compare como GCN e GAT atribuem peso à mensagem de cada vizinho durante a agregação. Qual afirmação descreve corretamente a diferença entre os dois mecanismos?
    \begin{enumerate}[(a)]
        \item No GCN o peso é fixado pelos graus dos nós, enquanto no GAT o peso é função dos embeddings dos nós envolvidos.
        \item No GCN o peso é aprendido por descida de gradiente, enquanto no GAT o peso é fixado pela estrutura do grafo.
        \item Ambos aprendem os pesos, mas o GAT restringe a soma dos pesos de cada nó a um.
        \item Ambos usam pesos fixos, e a diferença está apenas na função de ativação aplicada após a agregação.
        %a
    \end{enumerate}

    \item Uma equipe empilha $12$ camadas de GCN para classificação de nós em um grafo de citações de diâmetro pequeno e observa que a acurácia de validação cai em relação a um modelo de $3$ camadas. Qual fenômeno explica essa queda e qual intervenção o ataca diretamente?
    \begin{enumerate}[(a)]
        \item \textit{Overfitting} de parâmetros, mitigado por aumentar a taxa de aprendizado.
        \item \textit{Oversmoothing} dos embeddings, mitigado por \textit{skip connections} entre camadas.
        \item \textit{Vanishing gradient} na entrada, mitigado por remover a normalização de grau.
        \item Desbalanceamento de classes, mitigado por reamostrar os nós de treino.
        %b
    \end{enumerate}

    \item \emoji{skull-and-crossbones} Considere dois grafos não isomorfos, ambos com todos os nós de grau $2$ e atributos iniciais idênticos: $G_1$ é um ciclo de seis nós e $G_2$ são dois triângulos disjuntos. Sobre a capacidade de uma GNN de \textit{message passing} de distinguir esses dois grafos, qual afirmação é correta?
    \begin{enumerate}[(a)]
        \item Distingue após uma camada, pois os graus diferem entre os grafos.
        \item Distingue se a agregação usar soma em vez de média, pois a soma é injetiva.
        \item Não distingue, pois em toda iteração os \textit{multisets} de rótulos coincidem entre os grafos.
        \item Não distingue, pois ambos os grafos têm o mesmo número de arestas.
        %c
    \end{enumerate}

\end{enumerate}

\subsection*{Questões Dissertativas.}

\begin{enumerate}[I)]

    \item (\textbf{Propagação do GCN à mão}) Considere o caminho $G$ com $V = \{1,2,3\}$ e arestas $E = \{(1,2),(2,3)\}$, atributos escalares iniciais $\mathbf{h}^{(0)} = (2,\, 4,\, 6)^\top$, $W^{(1)} = 1$ e $\sigma = \mathrm{Id}$.
    \begin{enumerate}[(a)]
        \item Escreva $\tilde{\mathbf{A}} = \mathbf{A} + \mathbf{I}$ e o vetor de graus $\tilde{\mathbf{d}}$ no grafo com auto-laços.
        \item Calcule a matriz normalizada $\hat{\mathbf{A}} = \tilde{\mathbf{D}}^{-1/2} \tilde{\mathbf{A}} \tilde{\mathbf{D}}^{-1/2}$, entrada a entrada.
        \item Calcule $\mathbf{h}^{(1)} = \hat{\mathbf{A}}\, \mathbf{h}^{(0)}$ e interprete por que o valor do nó central difere dos valores das pontas.
    \end{enumerate}

    \item (\textbf{Tarefas e \textit{readout}}) Para cada uma das três tarefas canônicas em grafos, descreva o que se prevê e qual estratégia de \textit{readout} é adequada.
    \begin{enumerate}[(a)]
        \item \textit{Node classification}: o que é previsto e como o embedding por nó é usado na saída.
        \item \textit{Link prediction}: o que é previsto e como combinar os embeddings dos dois nós da aresta candidata.
        \item \textit{Graph classification}: o que é previsto e por que o \textit{readout} precisa ser invariante à permutação dos nós.
    \end{enumerate}

    \item \emoji{skull-and-crossbones} (\textbf{Soma, média e expressividade}) Considere um nó $v$ com dois vizinhos e um nó $v'$ com um único vizinho, todos os vizinhos carregando o mesmo atributo escalar $a$. Os \textit{multisets} de vizinhos são $\{\!\!\{a, a\}\!\!\}$ para $v$ e $\{\!\!\{a\}\!\!\}$ para $v'$.
    \begin{enumerate}[(a)]
        \item Calcule o resultado de agregar por média e por soma em cada caso e indique qual operação distingue $v$ de $v'$.
        \item Explique por que essa propriedade é descrita como injetividade sobre \textit{multisets} e relacione-a à camada do GIN, $\mathbf{h}^{(k)}_v = \mathrm{MLP}^{(k)}((1 + \epsilon^{(k)}) \mathbf{h}^{(k-1)}_v + \sum_{u \in \mathcal{N}(v)} \mathbf{h}^{(k-1)}_u)$.
        \item Dê um exemplo de tarefa em que essa distinção entre vizinhanças importa na prática e justifique.
    \end{enumerate}

    \item (\textbf{Profundidade e \textit{oversmoothing}}) Uma equipe quer aumentar o \textit{receptive field} de uma GCN para capturar dependências de longo alcance e cogita empilhar muitas camadas.
    \begin{enumerate}[(a)]
        \item Explique o que acontece com os embeddings dos nós à medida que o número de camadas cresce em um grafo de diâmetro pequeno, e por que isso degrada o desempenho.
        \item Descreva duas mitigações vistas em aula e explique, para cada uma, como ela ataca o problema.
        \item Comente o \textit{trade-off} entre profundidade, \textit{receptive field} e diversidade dos embeddings, indicando por que não é possível maximizar os três ao mesmo tempo.
    \end{enumerate}

\end{enumerate}

\newpage
\subsection*{Gabarito}
\color{blue}

\paragraph{Múltipla Escolha.}
\begin{enumerate}[I)]
    \item \textbf{(a)} Contando as arestas incidentes a cada nó: o nó $1$ aparece só em $(1,2)$, logo $d_1 = 1$; o nó $2$ aparece em $(1,2),(2,3),(2,4)$, logo $d_2 = 3$; o nó $3$ aparece em $(2,3),(3,4)$, logo $d_3 = 2$; o nó $4$ aparece em $(2,4),(3,4)$, logo $d_4 = 2$. O vetor $(2,4,3,3)$ corresponde aos graus no grafo com auto-laços $\tilde{\mathbf{A}} = \mathbf{A} + \mathbf{I}$, não ao grafo original. As outras opções não batem com a lista de arestas.

    \item \textbf{(b)} A equivariância à permutação exige que reordenar os vizinhos não altere o resultado da agregação, ou seja, que $\mathrm{AGG}$ seja invariante à ordem de seus argumentos (soma, média, máximo e atenção satisfazem isso). Linearidade não é necessária nem suficiente. Depender do grau ou ser injetiva sobre vetores individuais são propriedades distintas que não garantem invariância à ordem.

    \item \textbf{(a)} A vizinhança de $2$ é $\{1,3,4\}$, então $\mathcal{N}(2) \cup \{2\} = \{1,2,3,4\}$ com os quatro valores $(1,2,3,4)$. A média é $(1+2+3+4)/4 = 10/4 = 2{,}5$. A opção $3{,}0$ ignora o próprio nó na contagem, $4{,}5$ usa apenas alguns vizinhos e $2{,}0$ não corresponde a nenhuma agregação coerente.

    \item \textbf{(a)} O fator $\tilde{\mathbf{D}}^{-1/2} \tilde{\mathbf{A}} \tilde{\mathbf{D}}^{-1/2}$ depende apenas de $\tilde{\mathbf{A}}$ e $\tilde{\mathbf{D}}$, que são determinados pela estrutura do grafo e não contêm parâmetros treináveis, por isso é constante durante o treino. Ele não é produzido por \textit{backpropagation}, não depende de $\mathbf{W}^{(k)}$ e não incorpora $\sigma$, que é aplicada fora dele.

    \item \textbf{(b)} Com auto-laços, os graus são $\tilde d_2 = 4$ e $\tilde d_3 = 3$, então o peso é $1/\sqrt{4 \cdot 3} = 1/\sqrt{12} \approx 0{,}29$. O valor $0{,}41 = 1/\sqrt{6}$ usaria graus $2$ e $3$, $0{,}50$ ignoraria a normalização cruzada e $0{,}33$ corresponderia a uma média simples por grau.

    \item \textbf{(c)} Os nós $3$ e $4$ são estruturalmente simétricos: ambos têm grau $2$ e estão ligados ao nó $2$ e entre si, formando vizinhanças intercambiáveis. Com atributos iniciais iguais, o GCN não consegue diferenciá-los e produz o mesmo embedding. O nó $1$ tem grau $1$ e o nó $2$ tem grau $3$, então nenhum dos outros pares é simétrico.

    \item \textbf{(b)} GraphSAGE aprende uma função de agregação que age sobre vizinhanças locais, então pode ser aplicada a grafos novos sem reprocessar o grafo de treino: é indutivo por construção. O GCN na forma matricial precisa do grafo inteiro para montar o fator de propagação, o que o torna naturalmente transdutivo. A amostragem de vizinhos não dispensa os atributos de nó.

    \item \textbf{(c)} A $\mathrm{softmax}$ dos \textit{scores} $(2{,}0,\, 1{,}0,\, 1{,}5)$ dá $\exp(2{,}0) \approx 7{,}39$, $\exp(1{,}0) \approx 2{,}72$ e $\exp(1{,}5) \approx 4{,}48$, com soma $\approx 14{,}59$. Então $\alpha_{v,u_1} \approx 7{,}39/14{,}59 \approx 0{,}51$. Os valores $0{,}33$, $0{,}19$ e $0{,}67$ correspondem a outras normalizações ou aos demais coeficientes.

    \item \textbf{(a)} No GCN o peso de cada vizinho é $1/\sqrt{\tilde d_v \tilde d_u}$, fixado pelos graus e idêntico para qualquer tarefa. No GAT o coeficiente $\alpha_{vu}$ é obtido de $e_{vu} = \mathrm{LeakyReLU}(\mathbf{a}^\top [\mathbf{W} \mathbf{h}_v \Vert \mathbf{W} \mathbf{h}_u])$, ou seja, depende do conteúdo dos embeddings e se ajusta à tarefa. As demais opções trocam os papéis ou descrevem ambos como fixos.

    \item \textbf{(b)} Empilhar muitas camadas em um grafo de diâmetro pequeno leva ao \textit{oversmoothing}: após poucas camadas cada nó já vê quase todo o grafo e os embeddings convergem para vetores quase idênticos, perdendo poder discriminativo. \textit{Skip connections} atacam isso preservando a informação da camada anterior em cada passo. As outras opções diagnosticam fenômenos que não explicam a queda observada com mais camadas.

    \item \textbf{(c)} Ambos os grafos têm todos os nós de grau $2$, então cada nó sempre computa o mesmo $\mathrm{HASH}(c, \{\!\!\{c, c\}\!\!\})$ e os \textit{multisets} de rótulos coincidem em toda iteração do 1-WL. Como toda MPNN é no máximo tão poderosa quanto o 1-WL, nenhuma camada de \textit{message passing} separa $G_1$ de $G_2$. A soma não ajuda aqui, pois os \textit{multisets} de entrada já são idênticos; ter o mesmo número de arestas é consequência, não a causa.
\end{enumerate}

\paragraph{Dissertativas.}
\begin{enumerate}[I)]
    \item
    \begin{enumerate}[(a)]
        \item Para o caminho $1$, $2$, $3$, a matriz de adjacência é
        \[
            \mathbf{A} = \begin{bmatrix} 0 & 1 & 0 \\ 1 & 0 & 1 \\ 0 & 1 & 0 \end{bmatrix},
            \qquad
            \tilde{\mathbf{A}} = \mathbf{A} + \mathbf{I} = \begin{bmatrix} 1 & 1 & 0 \\ 1 & 1 & 1 \\ 0 & 1 & 1 \end{bmatrix}.
        \]
        Os graus com auto-laços são as somas das linhas: $\tilde{\mathbf{d}} = (2,\, 3,\, 2)$.
        \item Cada entrada é $\hat A_{ij} = \tilde A_{ij}/\sqrt{\tilde d_i \tilde d_j}$. Com $\tilde{\mathbf{d}} = (2,3,2)$:
        \[
            \hat{\mathbf{A}} = \begin{bmatrix}
                1/2 & 1/\sqrt{6} & 0 \\
                1/\sqrt{6} & 1/3 & 1/\sqrt{6} \\
                0 & 1/\sqrt{6} & 1/2
            \end{bmatrix}
            \approx \begin{bmatrix}
                0{,}50 & 0{,}41 & 0 \\
                0{,}41 & 0{,}33 & 0{,}41 \\
                0 & 0{,}41 & 0{,}50
            \end{bmatrix}.
        \]
        \item Aplicando a $\mathbf{h}^{(0)} = (2,4,6)^\top$:
        \begin{align*}
            h^{(1)}_1 &\approx 0{,}50 \cdot 2 + 0{,}41 \cdot 4 = 2{,}63, \\
            h^{(1)}_2 &\approx 0{,}41 \cdot 2 + 0{,}33 \cdot 4 + 0{,}41 \cdot 6 = 4{,}60, \\
            h^{(1)}_3 &\approx 0{,}41 \cdot 4 + 0{,}50 \cdot 6 = 4{,}63,
        \end{align*}
        ou seja, $\mathbf{h}^{(1)} \approx (2{,}63,\, 4{,}60,\, 4{,}63)^\top$. O nó central $2$ recebe contribuições dos três valores ($1$, $3$ e ele próprio), com peso menor sobre si mesmo por ter grau maior, enquanto as pontas $1$ e $3$ agregam apenas dois valores com peso $0{,}50$ sobre si, o que explica a diferença entre o valor central e os das pontas.
    \end{enumerate}

    \item
    \begin{enumerate}[(a)]
        \item \textit{Node classification}: prevê-se o rótulo de cada nó (por exemplo, a subárea de um artigo em uma rede de citações). O \textit{readout} usa diretamente o embedding $\mathbf{h}_v$ do nó, passado por um classificador.
        \item \textit{Link prediction}: prevê-se se uma aresta $(u,v)$ existe ou existirá (por exemplo, recomendar um item a um usuário). O \textit{readout} combina os embeddings $\mathbf{h}_u$ e $\mathbf{h}_v$, por concatenação, produto interno ou um MLP sobre o par, para produzir um \textit{score} de existência.
        \item \textit{Graph classification}: prevê-se um rótulo do grafo inteiro (por exemplo, se uma molécula é tóxica). O \textit{readout} agrega todos os $\mathbf{h}_v$ por \textit{pooling} (soma, média, máximo ou atenção) para obter $\mathbf{h}_G$. Esse \textit{pooling} precisa ser invariante à permutação porque o grafo não tem ordem canônica de nós: reordenar os nós não pode mudar a predição sobre o grafo.
    \end{enumerate}

    \item
    \begin{enumerate}[(a)]
        \item Média: para $v$, $\tfrac{a+a}{2} = a$; para $v'$, $\tfrac{a}{1} = a$. Os dois resultados coincidem, então a média não distingue. Soma: para $v$, $a + a = 2a$; para $v'$, $a$. Como $2a \neq a$ (para $a \neq 0$), a soma distingue $v$ de $v'$.
        \item A soma é injetiva sobre \textit{multisets} porque preserva a multiplicidade dos elementos: \textit{multisets} diferentes (como $\{\!\!\{a,a\}\!\!\}$ e $\{\!\!\{a\}\!\!\}$) produzem somas diferentes, ao passo que média e máximo descartam quantos elementos havia. A camada do GIN usa exatamente uma soma sobre os vizinhos seguida de um MLP, justamente para manter essa injetividade e atingir o teto de expressividade dado pelo teste 1-WL.
        \item Em classificação de moléculas, por exemplo, um átomo de carbono ligado a dois hidrogênios é quimicamente diferente de um ligado a um único hidrogênio. Se a agregação usar média ou máximo e os atributos dos vizinhos forem iguais, o modelo não vê a diferença de multiplicidade; com soma, vê. Por isso GIN costuma vencer GCN e GAT em tarefas de classificação de grafo que dependem de subestruturas.
    \end{enumerate}

    \item
    \begin{enumerate}[(a)]
        \item Cada camada faz o nó agregar informação de mais um salto, então após $K$ camadas o embedding de $v$ reúne informação a até $K$-hops. Em um grafo de diâmetro pequeno, depois de poucas camadas todos os nós já enxergam praticamente o grafo inteiro e suas representações passam a misturar as mesmas mensagens repetidamente, convergindo para embeddings quase idênticos. Esse é o \textit{oversmoothing}: os nós deixam de ser distinguíveis e a acurácia cai.
        \item Duas mitigações vistas em aula: (i) \textit{skip connections} no estilo $\mathbf{h}^{(k)}_v = \mathbf{h}^{(k-1)}_v + \mathrm{GNN}^{(k)}(\ldots)$, que preservam a informação da camada anterior e impedem que o embedding seja totalmente substituído pela média dos vizinhos; (ii) \textit{Jumping Knowledge} (JK-Net), que concatena ou faz \textit{max-pool} das ativações de todas as camadas no final, recuperando representações menos suavizadas das camadas iniciais. (Outras respostas válidas: \textit{DropEdge}, normalizações como \textit{PairNorm}, ou atenção global de Graph Transformers.)
        \item Aumentar a profundidade amplia o \textit{receptive field}, mas, passado o diâmetro do grafo, esse ganho vem acompanhado de \textit{oversmoothing}, que reduz a diversidade dos embeddings. Manter embeddings diversos favorece camadas rasas, que por sua vez limitam o \textit{receptive field}. Como aumentar a profundidade simultaneamente eleva o \textit{receptive field} e derruba a diversidade, não há como maximizar os três ao mesmo tempo: é preciso escolher um ponto de equilíbrio, em geral poucas camadas combinadas com mitigações.
    \end{enumerate}
\end{enumerate}

\color{black}

\end{document}
